% Chapter 7: Closing notes and pointers to v2

\chapter{Closing Notes and Forward Pointers}\label{chap:closing}

This document covers the core story of the kernel plus the
Borel-analytic separation that was discovered during formalization.
Three strands of the kernel are deferred to a successor blueprint.

\begin{itemize}[leftmargin=*,itemsep=4pt]
  \item \textbf{Compression.} The kernel formalizes a compression
    scheme for concept classes via an approximate minimax argument
    using multiplicative weights, giving a constructive route to
    compression schemes where the classical Moran--Yehudayoff
    argument~\cite{MoranYehudayoff2016} uses the exact minimax
    theorem. This material belongs to a dedicated v2 chapter that
    will import the textbook's Chapter~11 on compression and add the
    kernel-specific MWU formalization on top.

  \item \textbf{The measurable batch learner monad.} The kernel
    introduces \texttt{MeasurableBatchLearner}, a regularity axis
    isolating joint learner measurability, and proves that it forms a
    monad closed under Boolean combination, majority vote, piecewise
    interpolation, and countable selection. The first proof that
    non-neural learners (version spaces) satisfy
    \texttt{MeasurableBatchLearner} uses a rectangle decomposition
    that bypasses the Kuratowski--Ryll-Nardzewski measurable
    selection theorem for countable families. This material belongs
    to a dedicated v2 chapter on learner regularity.

  \item \textbf{PAC-Bayes, extended criteria, and the Baxter
    multi-task base case.} These are additional characterization
    results the kernel formalizes but which sit outside the core
    story of v1. They will be covered as separate strands in v2
    (PAC-Bayes as a frequentist--Bayesian bridge following
    McAllester~\cite{McAllester1999}, the extended criteria as a
    refinement of the measurability hierarchy, and the Baxter base
    case~\cite{Baxter2000} as a multi-task analogue of the
    fundamental theorem).
\end{itemize}

The kernel also contributes a proof engineering layer, the typed
proof operad \texttt{TPG\_FLT} together with the measurable inner
event metaprogram, that sits orthogonal to the mathematical content.
That layer is documented in a companion artifact at
\url{https://zetetic-dhruv.github.io/formal-learning-theory-kernel/methodology/}
(in preparation) which reuses the same leanblueprint infrastructure
and is cross-linked with this blueprint on the case study of the
measurable inner event pattern used by the symmetrization step of
Chapter~\ref{chap:measurability}.

For the full mathematical exposition of every strand deferred above,
together with the six paradigms not covered by the kernel (Ch~8
exact learning, Ch~9 universal learning, Ch~12 generalization
bounds, Ch~13 mind-change ordinals, Ch~15 analogies, Ch~17 extensions,
Ch~18 frontiers), see the companion textbook whose chapters this
blueprint reuses.
